\begin{salyga}
Su kuriomis parametro $m$ reikšmėmis lygtis turi vieną sprendinį?
    \[ x^2 - 2(m - 1)x + m^2 - 3 = 0 \]
\end{salyga}

\begin{atsakymas}
$m=2$.
\end{atsakymas}

\begin{sprendimas}
Kadangi $x^2$ koeficientas lygus $1$, lygtis turi vieną realųjį sprendinį tada ir tik tada, kai $D=0$. Čia
\[
\frac D4=(m-1)^2-(m^2-3)=4-2m.
\]
Taigi $D=0$ tada ir tik tada, kai $m=2$.
\end{sprendimas}
